% ch02-data — the data, the formal observation model, and the defect ledger

\section{Sources and window}\label{ch02:sec:sources}

Two data sources feed every model in this document.

\paragraph{PitchBook: the event-driven layer.} The deal history of the
private market, delivered as three relational tables whose column-level
dictionary is given in \cref{app:columns}. The \emph{deals} table is the
event log: one row per financing event, keyed by \code{deal\_id} and
\code{company\_id}, carrying the deal date, sector taxonomy
(\code{primary\_industry\_sector} / \code{group} / \code{code}), round
descriptors (\code{vc\_round}, \code{deal\_type}, \code{deal\_class},
up/down/flat flags), economics (\code{deal\_size} with a status flag
distinguishing \code{Actual}/\code{Estimated}/\code{Undisclosed},
\code{pre\_money\_valuation}, \code{post\_valuation},
\code{raised\_to\_date}, invested capital and ownership fields), firm
financials as reported around the deal (\code{revenue}, \code{ebitda},
\code{gross\_profit}, \code{net\_income}, employees, growth since last
deal), and geography. The \emph{companies} table is a snapshot per tracked
firm: identity, sector, founding year, business status, and
last-observed-deal summaries (\code{last\_financing\_date}, \code{\_size},
\code{\_valuation}, \code{\_deal\_type}). The \emph{investors} table links
firms to their investors with first-investment dates. The universe flag
\code{company\_universe} marks whether a deal belongs to a firm PitchBook
tracks; deals do occur with \code{company\_universe = 'Out of Universe'},
whose firms have \emph{no row} in the companies table --- defect (C6)
below.

\paragraph{FactSet: the market-clock layer.} A small panel of macro series
in FactSet Standardized Economics form: policy and effective rates
(\code{FFUNDT}, \code{FFUND}, daily), inflation (\code{CPI\_RAW} level and
\code{CPI\_YOY}, monthly), unemployment (\code{RUNEMP}, monthly), each
published with a lag of days to weeks after its reference period. These
series are the exogenous covariates of the ground process: they move on
calendar time regardless of deal activity, which makes them the only
covariates in the system that are \emph{never stale} in the sense of
\cref{ch:stale} --- but they are \emph{lagged}, and the lag matters for
point-in-time discipline (\cref{ch02:sec:pit}).

\paragraph{Window.} The joined history runs from January 2011 through June
2024: $T \approx 13.5$ years, i.e.\ roughly $704$ weeks of joint
deal-and-macro history. All development described in this document uses a
temporal split: a training window, a validation window for model selection,
and a strictly held-out test window at the end (the split protocol is fixed
in \cref{ch:gof}). The 2011 start postdates the crisis regime shift of
2008--2010 by design; the June 2024 end is an administrative censoring
time in the precise sense of defect (C1).

\section{The formal observation model}\label{ch02:sec:formal}

We now fix, once, the probabilistic description of what the data is. Every
later chapter states its assumptions relative to this model.

\begin{definition}[Market process]\label{def:ch02:market}
Work on a filtered probability space $(\Omega,\F,\{\Ft\},\Prob)$. The
market is a marked point process
\[
\{(t_m,\, s_m,\, i_m,\, \dmark_m)\}_{m \ge 1},
\qquad
t_1 \le t_2 \le \cdots,
\]
where $t_m \in [0,\Tend]$ is the (true) time of the $m$-th deal, $s_m \in
\sectors$ its sector, $i_m \in \firms \cup \firms^{\dagger}$ the receiving
firm --- $\firms$ the tracked universe, $\firms^{\dagger}$ the untracked
remainder --- and $\dmark_m$ the deal marks (size, valuation, round type,
\dots). For each sector $s$, $N_s(t)$ counts deals in $s$ up to $t$; for
each tracked firm $i$, $N_i(t)$ counts firm $i$'s deals. The filtration
$\Ft$ contains the internal history of the marked process and the paths of
all exogenous covariates up to $t$.
\end{definition}

\begin{definition}[Covariate processes]\label{def:ch02:covariates}
Three covariate processes enter the models:
\begin{enumerate}[label=(\roman*)]
\item \emph{Macro covariates} $\xcov^{\mathrm{mac}}(t) \in \R^{p}$:
exogenous, c\`adl\`ag, observed for all $t$ but only through their
\emph{published} version $\xcov^{\mathrm{mac}}_{\mathrm{pub}}(t) =
\xcov^{\mathrm{mac}}(t - L(t))$, where $L(t)$ is the publication lag
process (\cref{ch02:sec:pit}).
\item \emph{Firm covariates} $\xcov_i(t) \in \R^{q}$ for $i \in \firms$:
the true, latent, evolving state of firm $i$ (scale, financials,
capitalization). The data does \emph{not} deliver $\xcov_i(t)$; it
delivers the last value recorded at firm $i$'s most recent deal,
\begin{equation}\label{eq:ch02:locf}
\tilde{\xcov}_i(t) \;=\; \xcov_i\!\big(\tau_i(t)\big),
\qquad
\tau_i(t) \;=\; \sup\{t_m < t : i_m = i\},
\end{equation}
the last-observation-carried-forward (LOCF) covariate, whose observation
times are themselves event times of the modeled process. This is defect
(C5), the subject of \cref{ch:stale}.
\item \emph{Engineered history covariates}: deterministic functionals of
the observed history $\Hist_{t^-}$ chosen by us --- time since firm $i$'s
last deal, deal counts in trailing windows, sector activity summaries,
staleness ages. These are always available and always point-in-time
correct; \cref{ch:hawkes} uses them to build excitation \emph{and}
inhibition into intensity models, and \cref{ch:buckets} feeds them to
gradient boosting.
\end{enumerate}
\end{definition}

\begin{definition}[Conditional intensity and at-risk indicator]\label{def:ch02:intensity}
For a counting process $N_e$ ($e$ a sector, firm, or firm-transition
index), a \emph{conditional intensity} $\lam_e(t)$ is a nonnegative
predictable process such that
\begin{equation}\label{eq:ch02:martingale}
M_e(t) \;=\; N_e(t) - \int_0^{t} \atrisk_e(u)\, \lam_e(u) \dd u
\end{equation}
is a local martingale, where $\atrisk_e(t)\in\{0,1\}$ is the at-risk
indicator: the unit is under observation and eligible for the event at
$t^-$ \citep{abgk1993}. All likelihoods in this document are instances of
the point-process (Jacod) likelihood
\begin{equation}\label{eq:ch02:jacod}
\loglik(\thetamodel)
 \;=\;
 \sum_{m} \log \lam_{e_m}(t_m;\thetamodel)
 \;-\;
 \sum_{e} \int_{0}^{\Tend} \atrisk_e(u)\, \lam_e(u;\thetamodel) \dd u,
\end{equation}
whose martingale estimation theory (unbiased score at the truth, Fisher
information as predictable variation, asymptotic normality) is developed in
\cref{ch:foundations} following \citet{abgk1993} and
\citet{andersengill1982}.
\end{definition}

Most of the data's pathology concentrates in the at-risk indicator:
nearly every defect below corrupts either $\atrisk$ (who is really at
risk, and since when), or the event times $t_m$ (when did it really
happen), or the covariates the intensity conditions on.

\section{The defect ledger}\label{ch02:sec:defects}

Six defects, (C1)--(C6), exhaust the ways this dataset deviates from the
ideal of \cref{def:ch02:market}. For each: the formal statement, the
consequence if ignored, and the chapter that owns the treatment. The first
four restate, in final form, defects identified in the repository's
earlier theory notes; (C5) and (C6) are treated as first-class defects
here because the data review showed they bind hardest.

\subsection*{(C1) Administrative right censoring at the window end}

Every firm alive in June 2024 has its next round unobserved. Formally,
censoring is the predictable process $C_i(t) = \ind{t \le \Tend}$;
under independent censoring the observed process $\int C_i \dd N_i$ has
intensity $C_i(t)\lam_i(t)$ and the likelihood \cref{eq:ch02:jacod}
computed on observed data is a genuine partial likelihood
\citep[III.2]{abgk1993}. Administrative censoring at fixed $\Tend$ is
independent by construction. The censored spell contributes exactly its
survival factor $\exp(-\int \lam)$: no term dropped, no term invented.
This is the benign defect --- \emph{provided open spells are kept}.

\begin{warning}\label{warn:ch02:openspells}
Dropping open (censored) spells converts right censoring into right
\emph{truncation}: selection on the event occurring in-window, which
oversamples short inter-deal gaps and biases every fitted hazard upward.
This is not classical length bias (which oversamples \emph{long} spells in
prevalent-cohort designs); the correct treatment of sampling schemes and
truncation is \citet{kalbfleisch2002}. Policy: every dataset builder in
the repository must emit censored spells with an explicit censoring flag,
and every likelihood must consume them. \Cref{ch:poisson,ch:censoring}
implement this for Poisson and Hawkes calibration respectively;
\cref{ch:buckets} implements the bucketed analogue (censored buckets are
labeled, not silently dropped).
\end{warning}

\subsection*{(C2) Left truncation: firms enter observation by raising}

A firm is invisible to PitchBook until its first recorded financing: it
enters observation at $\entry_i$ = first-deal date, and enters
\emph{because of an event}. Mechanically this is delayed entry
\citep{keiding1990}: risk sets are $\atrisk_i(t) = \ind{\entry_i \le t <
U_i}$ and every integral in \cref{eq:ch02:jacod} starts at $\entry_i$.
Statistically it is a selection: the observed universe is ``firms
conditioned on having raised at least once,'' so unconditional statements
about companies are not identified from tracked data alone.
\Cref{ch:universe} is where this selection is modeled rather than ignored.

\subsection*{(C3) Unlabeled exits: the immortal-firm problem}

Failures and quiet wind-downs are mostly unrecorded (\code{business\_status}
is coarse and laggy). A dead firm keeps $\atrisk_i(t)=1$ in any naive risk
set. Three facts must be kept distinct. (i) The contamination bias ---
dead firms held at risk deflate every fitted hazard --- arises under
\emph{any} unlabeled-exit mechanism, even exit independent of funding
propensity: it is a misrecorded $\atrisk_i$, not a dependence effect.
(ii) Because exit also correlates with the latent quality that drives
funding, we do not expect the defect to be correctable from observables:
even a good exit classifier is unlikely to reconstruct truthful risk sets,
and we expect risk-set truthfulness to be one of the largest sensitivities
in the whole system. Both expectations are to be tested on the synthetic
world (\cref{ch:results-synth}). (iii) If
exits \emph{were} labeled, removing firms at exit would be valid for
cause-specific funding intensities with no independence assumption at all
--- the standard competing-risks fact \citep{putter2007}. The problem is
that exits are unobserved, not that they are dependent. Treatments ---
exposure windows, mover--stayer mixtures \citep{farewell1982}, and
ranking-layer mitigations --- are developed where they bind:
\cref{ch:poisson,ch:buckets,ch:ranking}.

\subsection*{(C4) Untrusted deal dates and aggregation}

Two distinct corruptions of $t_m$. First, \emph{date jitter}: PitchBook
records announced vs.\ closed dates that can differ by days, and a
nontrivial fraction of dates are back-filled or rounded; the recorded
$\hat t_m$ is best modeled as $t_m$ observed through a small-window error.
Second, \emph{deliberate aggregation}: for modeling at weekly resolution we
replace times by bin counts, which is interval censoring by construction.
Both are instances of one statement: the event count over an interval is
trusted, the position within the interval is not. \Cref{ch:censoring}
builds the corresponding likelihoods: interval-censored NHPP calibration
in closed form, and interval-censored Hawkes via the mean-intensity
(mean-Poisson-process) approximation of \citet{rizoiu2022}, with the
information loss quantified --- branching-scale parameters survive
aggregation, timescale parameters do not (the $\branch$--$\decay$ ridge of
\cref{ch:gof}; Fisher information in the decay $\to 0$ as bin width grows
past the kernel half-life).

\subsection*{(C5) Stale firm covariates, updated only at own deals}

Defect \cref{eq:ch02:locf}: firm covariates are observed only at the
firm's own deals, so at any $t$ the available $\tilde\xcov_i(t)$ has age
$\staleage_i(t) = t - \tau_i(t)$, and the age is largest exactly for the
firms that have gone longest without raising --- the observation times are
\emph{informative}, being event times of the process under study. This is
simultaneously a measurement-error problem (attenuation of covariate
effects; \citealp{carroll2006,prentice1982}) and an
internal-covariate problem in the sense of \citet{kalbfleisch2002}.
\Cref{ch:stale} formalizes both faces and commits to a layered remedy:
staleness-age features always; an explicit covariate-evolution state layer
where it pays; joint-model escalation
\citep{rizopoulos2012,tsiatisdavidian2004} only if a trigger fires. The
size of the attenuation is to be measured on the synthetic world
(\cref{ch:results-synth}).

\subsection*{(C6) Out-of-universe deals: events with no covariates}

Deals occur from firms not in the tracked universe
(\code{company\_universe = 'Out of Universe'}, no companies-table row, no
covariates, no history). Consequences: sector deal counts are complete but
firm-attributable counts are thinned; any firm-level risk set is missing
members who \emph{do} compete for deals; and the mark factor must own an
outside option or it will overstate every tracked firm's probability.
\Cref{ch:universe} develops the treatment: an explicit newcomer/outside
component in the mark factor (with the discipline that \emph{newcomer
probability} --- first appearance in the data --- must never be conflated
with \emph{entrant probability} --- economic entry), a ghost/external
component in the ground process, and unseen-mass estimation in the
Good--Turing spirit \citep{good1953,galesampson1995} as the model-free
check on how much deal mass the tracked universe can see.

\begin{table}[htbp]
\centering
\small
\begin{tabular}{p{0.06\textwidth} p{0.30\textwidth} p{0.28\textwidth} p{0.22\textwidth}}
\toprule
\textbf{Defect} & \textbf{What is corrupted} & \textbf{If ignored} &
\textbf{Owned by} \\
\midrule
(C1) & spells open at June 2024 & upward-biased hazards if open spells
dropped & \cref{ch:poisson,ch:censoring,ch:buckets} \\
(C2) & entry into observation ($\atrisk_i$ before first deal) &
selection mistaken for dynamics & \cref{ch:universe} \\
(C3) & $\atrisk_i$ after unrecorded death & deflated hazards, inverted
rankings & \cref{ch:poisson,ch:buckets,ch:ranking} \\
(C4) & event times $t_m$ & spurious timing estimates; broken
$\decay$ & \cref{ch:censoring,ch:gof} \\
(C5) & firm covariates $\xcov_i(t)$ & attenuated, staleness-confounded
effects & \cref{ch:stale} \\
(C6) & firm identity $i_m$, risk-set membership & overstated tracked-firm
probabilities & \cref{ch:universe} \\
\bottomrule
\end{tabular}
\caption{The defect ledger. Every modeling chapter states which defects it
is robust to and which assumptions it needs about the rest.}
\label{tab:ch02:defects}
\end{table}

\section{Point-in-time discipline}\label{ch02:sec:pit}

A model that consumes information not available at prediction time is
worthless regardless of its metrics. The repository already enforces the
following discipline, and this document makes it policy (implemented in
\code{synthetic/loaders.py}, to be preserved verbatim by every real-data
loader):

\begin{protocol}[Point-in-time construction]\label{prot:ch02:pit}
\leavevmode
\begin{enumerate}[label=(\alph*)]
\item \emph{Macro covariates are as-of joined on publish date}: the value
of $\xcov^{\mathrm{mac}}_{\mathrm{pub}}(t)$ is the last value whose
\emph{publication} date (not reference date) precedes $t$.
\item \emph{Firm features are LOCF from past deals only}: $\tilde\xcov_i(t)$
uses deals strictly before $t$; the staleness age $\staleage_i(t)$ is
always carried alongside as a feature, never discarded.
\item \emph{Risk-set membership is inferred from observables}: founding
year and first-deal date determine entry; any exit proxy used to close
spells must itself be point-in-time.
\item \emph{Deduplication and hygiene run before everything}: duplicate
deal rows (same firm within a few days, jittered size) are collapsed by
the documented rule; sector mislabels and undisclosed sizes are handled by
the imputation rules of \cref{app:columns}, never ad hoc.
\item \emph{Evaluation splits are temporal}, with model selection strictly
inside the training era. No hyperparameter, feature transform, or
calibration map may see post-split data.
\end{enumerate}
\end{protocol}

\section{The synthetic world}\label{ch02:sec:synthetic}

The real data's defects cannot be switched off, so the repository
maintains a synthetic world for validation: generated PitchBook-style
tables (deals/companies/investors) and a FactSet-style macro feed with
publication lags, passed through an observation layer that reproduces
defects (C1)--(C6) with known ground truth (\code{synthetic/}). Two
regimes matter: regime A is generated by the exact two-stage model family
we fit (a well-specified test), and regime B is a daily firm-level
block-Hawkes world with self-inhibition and latent quality, deliberately
\emph{not} in the fitted family (a misspecification stress test). The
synthetic world is a planned validation tool, not a source of standing
results. The experiments it supports are specified in
\cref{ch:results-synth}, under a fixed transfer convention: synthetic
outcomes set \emph{priors and vetoes} for real-data work, never
substitutes for real-data evaluation.

\begin{decision}{D2.1 --- defects are assumptions, stated per model}
Every model chapter opens by declaring, defect by defect against
\cref{tab:ch02:defects}, whether the model (i) is robust to the defect,
(ii) corrects it under stated assumptions, or (iii) is vulnerable and
accepts the bias with a documented sign. No silent vulnerability is
permitted. Rationale: we expect the largest errors to come not from wrong
models but from unstated exposure to (C3) and (C5); the synthetic
experiments of \cref{ch:results-synth} are designed to test this.
\end{decision}

\begin{codehooks}
Real-data ingestion will mirror \code{synthetic/loaders.py}: an adapter
\modrepo{data/pitchbook.py} (to be written) emitting the same
fitter-ready arrays --- weekly sector counts, as-of macro matrix, LOCF firm
features \emph{plus staleness ages}, at-risk masks, and the event list ---
so that every fitter downstream is source-agnostic. The defect ledger of
\cref{tab:ch02:defects} becomes a validation suite
(\modrepo{data/audits.py}): open-spell retention check, entry-date
consistency, duplicate-collapse idempotence, as-of-join look-ahead test,
and universe-share reporting (\code{off\_universe\_share} is already a
first-class output of the synthetic loader and must be one for the real
loader).
\end{codehooks}
